All graphs considered in this paper are simple, i.e., we disallow multiple edges and self-loops. If $G=(V,E)$ is a directed graph, we call $S\subseteq V$ a \emph{closed} set if for all $(u,v)\in E$, we have that $u\in S$ implies that $v\in S$. We say that $v$ is a \emph{descendant} of $u$ if $v$ is reachable by a path from $u$ and let $D(u)$ denote the set of descendants of $u$ (including $u$ itself). Analogously, if $v$ is a descendant of $u$, we call $u$ an \emph{ascendant} of $v$. We extend the notation naturally to sets $S\subseteq V$ by defining $D(S) \coloneqq \bigcup_{u\in S} D(u)$.
For a graph $G=(V,E)$ with node weights $w: V\to \nats$ and $S\subseteq V$, we write $w(S) \coloneqq \sum_{v\in S}w(v)$. For any $S\subseteq V$, we let $G[S]$ denote the subgraph of $G$ induced by $S$, i.e., the subgraph obtained by deleting all vertices in $V\setminus S$ and adjacent edges.
%
